Figure~\ref{fig:landscape} shows why the grid matters so much. Within the training distribution every
kernel scores between $0.95$ and $0.99$ at every bandwidth for $C=1$, and that is also what the kernels
achieve on held-out rows of the training file: there, standard cross-validation ranks the 40 settings
correctly for every kernel ($r=+0.94$ to $+0.99$; Table~\ref{tab:decomp}) and the narrow-grid artefact is
absent ($\Delta$ROC-AUC $-0.002$). On the official test set the raw-feature RBF kernel is instead U-shaped
in $\gamma$: strong at very wide and very narrow bandwidths (up to $0.945$), and weakest at intermediate
ones ($0.740$ at worst across the grid), which is exactly where a grid capped at $\gamma=3$ leads
cross-validation. Across the 40 settings its standard CV score is uncorrelated with its test ROC-AUC
($r=-0.16$). The training rows and the cross-validation are the same in both cases; only the test
distribution differs, and a criterion computed on the training distribution cannot anticipate how the
shift reshapes the test landscape.

Table~\ref{tab:decomp} separates the shift into its parts. Removing the 17 novel attack types from the
test set (3{,}750 rows) leaves the blindness in place (RBF $r=-0.01$) and shrinks the narrow-grid
artefact from $0.167$ to $0.063$ ROC-AUC. Reweighting the remaining rows to the training file's
attack-type mix shrinks it to $0.012$, as the RBF landscape flattens (test ROC-AUC $0.97$ to $0.99$ at
every setting, which also makes its $r$ uninformative there). Of the aggregate artefact, the novel types
thus account for about $60\%$ and the changed attack mix for about $30\%$. At the operating point the
decomposition differs: the RBF setting selected on the narrow grid detects $0.078$ of attacks at $1\%$
false positives on the official set, $0.090$ without the novel types and $0.087$ after reweighting,
against $0.964$ on held-out training rows. That collapse comes from drift within attack types the
training file contains, not from the novel types.

@@TAB_DECOMP@@
